% from ega1-1-fr.tex
\oldpage[I]{90}
Si $M$, $N$ sont deux $B$-modules, le $\widetilde{B}$-Module
$\widetilde{M}\otimes_{\widetilde{B}}\widetilde{N}$ s'identifie
canoniquement à $\widetilde{(M\otimes_B N)}$~; de même
$\mathscr{H}\!om_{\widetilde{B}}(\widetilde{M},\widetilde{N})$ s'identifie
canoniquement à $\widetilde{(\operatorname{Hom}_B(M,N))}$ lorsque $M$ admet
une présentation finie~; les démonstrations sont les mêmes que dans
(\hyperref[I.1.3.12-fr]{1.3.12}).

Si $\mathfrak{J}$ est un idéal de $B$, $N$ un $B$-module, on a
$\widetilde{(\mathfrak{J}N)}=\widetilde{\mathfrak{J}}\cdot\widetilde{N}$.

Enfin, si $B$ est une $A$-algèbre graduée par des sous-$A$-modules $B_n$
($n\in\mathbb{Z}$), la $\widetilde{A}$-Algèbre $\widetilde{B}$, somme directe
des $\widetilde{A}$-Modules $\widetilde{B_n}$
(\hyperref[I.1.3.9-fr]{1.3.9}) est graduée par ces sous-$\widetilde{A}$-Modules,
l'axiome de la graduation exprimant que l'image de l'homomorphisme
$B_m\otimes_A B_n\to B$ est contenu dans $B_{m+n}$. De même, si $M$ est un
$B$-module gradué par des sous-modules $M_n$, $\widetilde{M}$ est un
$\widetilde{B}$-Module gradué par les $\widetilde{M_n}$.
\end{env}

\begin{env}[1.3.14]
\label{I.1.3.14-fr}
Si $B$ est une $A$-algèbre, $M$ un sous-module de $B$, la
sous-$\widetilde{A}$-Algèbre de $\widetilde{B}$ engendrée par $\widetilde{M}$
(\hyperref[0.4.1.3-fr]{0, 4.1.3}) est la sous-$\widetilde{A}$-Algèbre
$\widetilde{C}$, en désignant par $C$ la sous-algèbre de $B$ engendrée par
$M$. En effet, $C$ est la somme des sous-modules de $B$ images des
homomorphismes $\bigotimes_A^n M\to B$ ($n\geq0$), et il suffit d'appliquer
(\hyperref[I.1.3.9-fr]{1.3.9}) et (\hyperref[I.1.3.12-fr]{1.3.12}).
\end{env}

\subsection{Faisceaux quasi-cohérents sur un spectre premier.}
\label{subsection:I.1.4-fr}

\begin{theorem}[1.4.1]
\label{I.1.4.1-fr}
Soient $X$ le spectre premier d'un anneau $A$, $V$ une partie ouverte
quasi-compacte de $X$, $\mathscr{F}$ un $(\mathscr{O}_X|V)$-Module. Les
quatre conditions suivantes sont équivalentes~:
\begin{enumerate}
  \item[\textit{a})] Il existe un $A$-module $M$ tel que $\mathscr{F}$ soit
  isomorphe à $\widetilde{M}|V$.
  \item[\textit{b})] Il existe un recouvrement ouvert fini $(V_i)$ de $V$ par des
  ensembles de la forme $D(f_i)$ ($f_i\in A$) contenus dans $V$, tels que,
  pour tout $i$, $\mathscr{F}|V_i$ soit isomorphe à un faisceau de la forme
  $\widetilde{M_i}$, où $M_i$ est un $A_{f_i}$-module.
  \item[\textit{c})] Le faisceau $\mathscr{F}$ est quasi-cohérent
  (\hyperref[0.5.1.3-fr]{0, 5.1.3}).
  \item[\textit{d})] Les deux propriétés suivantes ont lieu~:
  \begin{enumerate}
    \item[\textit{d}\,1)] Pour tout $f\in A$ tel que $D(f)\subset V$ et toute
    section $s\in\Gamma(D(f),\mathscr{F})$, il existe un entier $n\geq0$ tel
    que $f^n s$ se prolonge en une section de $\mathscr{F}$ sur $V$.
    \item[\textit{d}\,2)] Pour tout $f\in A$ tel que $D(f)\subset V$ et toute
    section $t\in\Gamma(V,\mathscr{F})$ telle que la restriction de $t$ à
    $D(f)$ soit $0$, il existe un entier $n\geq0$ tel que $f^n t=0$.
  \end{enumerate}
\end{enumerate}
\end{theorem}

(Dans l'écriture des conditions \textit{d}\,1) et \textit{d}\,2), on a tacitement identifié $A$
et $\Gamma(\widetilde{A})$ en vertu de
\hyperref[I.1.3.7-fr]{1.3.7}.)

Le fait que \textit{a}) entraîne \textit{b}) est conséquence immédiate de
(\hyperref[I.1.3.6-fr]{1.3.6}) et du fait que les $D(f_i)$ forment une base
de la topologie de $X$ (\hyperref[I.1.1.10-fr]{1.1.10}). Comme tout
$A$-module est isomorphe au conoyau d'un homomorphisme
$A^{(I)}\to A^{(J)}$, (\hyperref[I.1.3.9-fr]{1.3.9}) prouve que tout faisceau
associé à un $A$-module est quasi-cohérent~; donc \textit{b}) entraîne \textit{c}).
Réciproquement, si $\mathscr{F}$ est quasi-cohérent, tout $x\in V$ possède
un voisinage de la forme $D(f)\subset V$ tel que $\mathscr{F}|D(f)$ soit
isomorphe au conoyau d'un homomorphisme
$\widetilde{A_f}^{(I)}\to\widetilde{A_f}^{(J)}$, donc au faisceau associé au
module $\widetilde{N}$, conoyau de l'homomorphisme
$A_f^{(I)}\to A_f^{(J)}$ correspondant
(\hyperref[I.1.3.8-fr]{1.3.8} et \hyperref[I.1.3.9-fr]{1.3.9})~; comme $V$
est quasi-compact, il est clair que \textit{c}) entraîne \textit{b}).
